\section{Introduction}

Axonal degeneration is a central pathological feature of neurodegenerative disease and a principal substrate of irreversible clinical disability in multiple sclerosis (MS) \cite{criste2014axonal,haines2011axonal}. In MS, axonal damage arises both as a direct consequence of inflammatory attack and secondarily to demyelination, glial activation and exposure to excitatory amino acids and cytokines, and the transition to progressive disease is thought to occur when axonal loss exceeds the compensatory capacity of the central nervous system \cite{criste2014axonal}. Detecting axonal pathology in vivo, early in the disease and noninvasively, is therefore a priority for improving prognosis, monitoring treatment response and guiding the development of neuroprotective therapies.

Magnetic resonance imaging (MRI), and in particular diffusion-weighted MRI, has been the workhorse of in vivo investigation of white matter microstructure in MS and related conditions \cite{inglese2010diffusion}. Conventional diffusion tensor imaging (DTI) is however notoriously nonspecific: measures such as fractional anisotropy conflate contributions from axons, myelin, extracellular space and glia, and cannot unambiguously localise axonal damage \cite{desantis2014accuracy}. Multicompartment models offer a principled alternative by partitioning the diffusion signal into geometric compartments whose parameters map onto biological quantities, and AxCaliber in particular provides an estimate of the mean axon diameter \cite{assaf2008axcaliber,alexander2010orientationally,sepehrband2016towards}. Because short diffusion times and very high gradient amplitudes are needed to sensitise the signal to micron-scale restriction, AxCaliber has historically been confined to small-bore preclinical scanners or to single-tract regions of interest on human systems \cite{jones2018microstructural,barazany2009invivo}. A recent clinical application reported an increase in axonal diameter in the normal-appearing white matter (NAWM) of the corpus callosum of MS patients \cite{desantis2019whole}, but in its original formulation AxCaliber is only applicable to voxels that contain a single dominant fiber orientation. Since at least 70\% of brain voxels contain two or more dominant fiber orientations \cite{jeurissen2013investigating}, whole-brain axonal-diameter mapping has remained out of reach.

Equally important, the interpretation of any imaging-based axonal-damage biomarker rests on histological validation, yet quantitative comparison between in vivo diffusion MRI and ex vivo tissue is rarely performed and is hampered by fixation-related distortions \cite{horowitz2015axon}. Previous studies have shown that AxCaliber-derived axon diameters correlate with electron-microscopy measurements in healthy rodent and post-mortem human tissue \cite{barazany2009invivo,horowitz2015axon}, but no work to date has validated AxCaliber against a controlled model of acute axonal damage. Existing preclinical MS models such as experimental autoimmune encephalomyelitis (EAE) and cuprizone-induced demyelination \cite{constantinescu2011eae,torkildsen2008cuprizone} primarily perturb the myelin or immune compartments rather than selectively the axon, leaving open the question of whether a measured increase in axon diameter truly reflects axonal pathology.

Here we close this gap with a translational pipeline. First, we validate an AxCaliber formulation modified to handle multiple fiber orientations \cite{desantis2019whole,tax2014recursive} against a preclinical model of acute axonal damage, in which unilateral stereotaxic injection of the NMDA agonist ibotenic acid into the rat hippocampus produces Wallerian-like degeneration of the fimbria while leaving the myelin sheath initially preserved \cite{conforti2014wallerian,adalbert2012intra}. Second, we apply a species-matched protocol on a 3T Connectom system equipped with 300\,mT/m gradients and a 64-channel array \cite{jones2018microstructural} to a cohort of relapsing-remitting MS patients (diagnosed according to the 2010 McDonald criteria \cite{polman2011diagnostic}) and age-matched healthy controls, quantifying disability with the Expanded Disability Status Scale (EDSS) \cite{kurtzke1983rating}.

Our contributions are: (i) the first preclinical validation of AxCaliber against histology in a model of acute axonal damage, demonstrating that a significant paired increase in estimated axonal diameter along the ibotenic-acid-injected fimbria ($p=0.003$) accompanies higher neurofilament intensity, lower NeuN staining and no change in myelin basic protein, with a significant Pearson correlation ($r=0.54$, $p=0.029$) between neurofilament fluorescence and MRI-derived diameter; (ii) the first whole-brain, crossing-fiber-compatible map of axonal-diameter changes in MS NAWM, showing a symmetric increase across every major tract, extending prior corpus-callosum findings \cite{desantis2019whole,ercan2015multimodal}; (iii) a lesion-stratification analysis showing that increased axonal diameter is selective for T1-hypointense lesions \cite{vanwaesberghe1999axonal}, consistent with their greater axonal burden; and (iv) a fully documented translational protocol that matches acquisition parameters between preclinical and clinical scanners, providing a template for future axon-centric imaging studies.

\section{Related Work}

\paragraph{Multicompartment diffusion models of axon diameter.}
AxCaliber \cite{assaf2008axcaliber} parameterises the diffusion-weighted signal as a sum of restricted intra-axonal and hindered extra-axonal components, from which a mean axon diameter can be recovered given sufficient $q$-space sampling and gradient strength. Orientationally invariant extensions such as ActiveAx \cite{alexander2010orientationally} relaxed the requirement for strictly parallel fibers but still struggle in voxels with multiple fiber populations; subsequent work has quantified and partly mitigated the resulting biases \cite{sepehrband2016towards,desantis2019whole}. Our implementation adopts a fiber-dispersion formulation \cite{desantis2019whole,tax2014recursive} that accounts for more than one fiber orientation per voxel, a necessary condition for whole-brain mapping since the majority of brain white-matter voxels contain crossing fibers \cite{jeurissen2013investigating}.

\paragraph{High-gradient clinical diffusion MRI.}
Measuring micron-scale axonal restriction with diffusion MRI requires short diffusion times and strong gradients, for which purpose the human Connectom platform (300\,mT/m) was developed \cite{jones2018microstructural}. Paired with a 64-channel array coil, this system allows $b$-values of up to 4000\,s/mm$^2$ with modest TE penalties, which we exploit in our clinical acquisition. The matched acquisition protocol (Table~4) preserves the key $b$-value and diffusion-time ranges between the 7T preclinical scanner and the 3T Connectom, enabling a like-for-like translational comparison \cite{desantis2019whole}.

\paragraph{Histological validation of diffusion MRI biomarkers.}
Quantitative validation of diffusion-MRI biomarkers against histology is comparatively rare and subject to well-known caveats: tissue fixation, shrinkage and limited field of view all distort the comparison \cite{horowitz2015axon}. Prior validation of AxCaliber \cite{barazany2009invivo,horowitz2015axon} established correlations with electron microscopy in healthy tissue and demonstrated that AxCaliber-derived diameters scale with conduction velocity, but none tested whether AxCaliber is sensitive to pathological alterations of the axonal compartment. Selective axonal injury via stereotaxic injection of neurotoxins is a well-established strategy in rodent neuroscience \cite{conforti2014wallerian,adalbert2012intra} and has recently been leveraged to dissect astrocytic and microglial contributions to gray-matter diffusion contrasts in a related paradigm, providing a precedent for the approach adopted here.

\paragraph{Axonal pathology in MS: biomarkers and lesion heterogeneity.}
Post-mortem and animal evidence converge on varicosities, spheroids and transport failure as the morphological hallmarks of axonal injury in MS, with smaller axons being more vulnerable \cite{criste2014axonal,tallantyre2010demyelinated,nikic2011reversible,bergers2002axonal}. Recent work shows that axon-myelin unit blistering can occur in NAWM even in the absence of overt demyelination or inflammation \cite{luchicchi2021axonal,teo2021nilered}. Fluid-phase neurofilament assays are now accepted surrogate markers of axonal injury \cite{petzold2005neurofilament}, and T1-hypointense (``black hole'') MS lesions are understood to harbour greater axonal loss than isointense lesions \cite{vanwaesberghe1999axonal,raz2015relationship}. Our results connect these two strands: a whole-brain, histologically validated imaging biomarker of axonal caliber that is elevated throughout NAWM and selectively elevated in T1-hypointense lesions.

\paragraph{Statistical frameworks for voxelwise white-matter analysis.}
For whole-brain voxelwise inference we build on tract-based spatial statistics (TBSS) \cite{smith2006tbss}, which projects per-subject diffusion measures onto a common white-matter skeleton to mitigate misregistration. Because nonlinear registration accuracy is critical at lesion boundaries, we adopt a high-accuracy diffeomorphic alignment \cite{klein2009evaluation} and perform lesion-masked permutation-based inference on the general linear model with threshold-free cluster enhancement \cite{winkler2014permutation}, controlling for age and for multiple comparisons across clusters.
